\documentclass[11pt]{article}
\usepackage[a4paper,margin=24mm]{geometry}
\usepackage{amsmath,amssymb,amsthm,mathtools,booktabs,microtype}
\usepackage[T1]{fontenc}
\usepackage[hidelinks]{hyperref}
\newtheorem{theorem}{Theorem}
\newtheorem{lemma}{Lemma}
\newtheorem{proposition}{Proposition}
\newcommand{\R}{\mathbb R}
\newcommand{\Q}{\mathbb Q}
\newcommand{\dd}{\mathrm d}
\newcommand{\one}{\mathbf 1}
\DeclareMathOperator{\diag}{diag}
\DeclareMathOperator{\tr}{tr}
\title{Exact C1 certificates in the order-13744 candidate}
\author{}
\date{23 October 2026}
\begin{document}
\maketitle
\begin{abstract}
An exact rational certificate places the previously separated C1 support
inside the canonical cone used by the order-13744 candidate.  A rational
selector extends this certificate to an open five-parameter contact family:
its compressed Grams have margin $2^{-22}$ at the reference source and
remain positive for target variation below $1/3284140032$.
A displayed one-parameter subfamily is certified on
$|\sigma-140|\leq2^{-42}$.  The same selector fails at other regular
sources, so universal adequacy remains unresolved.
\end{abstract}

The canonical C1 cone used by the order-13744 outer approximation contains
the physical support that separates the two smaller candidates.  More
generally, an explicit rational selector proves membership for an open
five-parameter family of stationary sources around that support.  The
selector does not work on the entire regular contact domain.  Neither a
universal converse for order 13744 nor a converse for C160 is proved here.

\section{The canonical polynomial and its finite Gram representation}

Write $c=(d,m,C,A,D,\alpha,\beta)$ and $\psi=(W,X,Y,Z)^T$.
Set $D_u=1+u^2$, $D_v=1+v^2$, and define the entire scalar parent by
\begin{align}
 F_c={}&4m(vW-X+uvY-uZ)^2+4CD_v(uW-Y)^2
             +4AD_uD_vY^2 \notag\\
 &+(4D+2d)(uW+uvX-Y-vZ)^2
 +(4\alpha+2d)D_u(vY-Z)^2 \notag\\
 &+(4\beta+2d)D_v(X+uZ)^2-2dD_uD_vZ^2.
 \label{eq:parent}
\end{align}
The symbol $\mathcal P_c$ denotes the symmetric matrix such that
$F_c=\psi^T\mathcal P_c\psi$.  All source coordinates are fixed before
the independent testing variables $u,v,\psi$ are introduced.

For each of the four weights
\[
 (\alpha_g,\beta_g)=(0,0),(0,1),(1,0),(1,1),
 \qquad g=0,1,2,3,
\]
start with the ordered features $u^rv^s\psi_i$, ordered lexicographically
by $(i,r,s)\in\{0,1,2,3\}\times\{0,1\}^2$.
Remove the following features:
\[
\begin{array}{c|c|c}
 g&\text{removed features}&n_g\\\hline
 0&W,\ uX&14\\
 1&W&15\\
 2&W,\ uX&14\\
 3&uX&15
\end{array}
\]
Let $V_g$ be the retained feature vector.  Define $K_1$ by
$d\geq0$ and the existence of $G_g\succeq0$ satisfying the polynomial
identity
\begin{equation}
 (1+u)(1+v)F_c
  =\sum_{g=0}^3u^{\alpha_g}v^{\beta_g}V_g^TG_gV_g.
 \label{eq:gram}
\end{equation}
There are 450 upper-triangular entries and 148 coefficient equations.
Each Gram entry contributes to one scalar monomial, so the coefficient
equations have disjoint supports.  Solving for one entry in each support
leaves 302 free entries.  These are the actual auxiliary coordinates in
\path{C1_ENLARGED_PRIMAL.coo.json}.  The four Gram blocks and $[d]$
have total order 59.  The global compiler additionally imposes the six
event signs and retains both canonical C1 orientations.

Equation \eqref{eq:gram} proves $F_c\geq0$ for every $u,v\geq0$ and every
state vector.  Its multiplier is strictly positive there.  The exact sign
gauges $(u,Y,Z)\mapsto(-u,-Y,-Z)$ and
$(v,X,Z)\mapsto(-v,-X,-Z)$ extend this conclusion to signed finite angles.
Continuity of the normalized physical operator gives the projective
endpoints.  Thus membership in $K_1$ is an all-state certificate, rather
than a test restricted to its source contact.

\section{An exact certificate for the separated support}

The source is
\begin{equation}
\begin{aligned}
 c_*={}&\left(1,
 \frac{1166523000}{31661820415017653},
 \frac{322803000}{186199784281},
 \frac{218288946200}{1359122513},\right.\\[-1mm]
 &\left.\hspace{13mm}
 \frac{90943434293}{222507771414},
 \frac{248802430236}{3300726103},
 \frac{205476973}{51646655494}\right),\\
 &\theta_*=(1/13,16,140,-1/32,1/50).
\end{aligned}
\label{eq:seed}
\end{equation}
Its homogeneous contact vector is $(140,-1/32,1/50,1)$, with last
coordinate fixed to one.  It is not a unit vector; dividing by its
Euclidean norm gives the corresponding pure quantum state.
The data file
\path{C1_REPAIRED_TWO_FRAME_COUNTEREXAMPLE.json}
supplies all 302 rational auxiliary entries $y_*$, all four actual Gram
matrices, and rational factorizations
\begin{equation}
 G_g(c_*,y_*)=F_g\Lambda_gF_g^T,
 \qquad \Lambda_g\succ0,
 \qquad \operatorname{rank}F_g=n_g-1.
 \label{eq:seedfact}
\end{equation}
The first $n_g-1$ rows of each $F_g$ are lower triangular with nonzero
diagonal.  Substitution into the current source pencil gives these exact
matrices.  Their quadratic forms give \eqref{eq:gram}, and the four state
kernel equations and both fixed-state angular derivatives vanish at
\eqref{eq:seed}.  These are rational identities and positive rational
diagonal entries; no optimization status is used.

The canonical source pencil is byte-identical to the source array used
by the current order-13744 construction; its SHA256 is
\begin{center}\small\ttfamily
1e1eb340077940be9c466f176efc4816d95c5271d448b7b52e94aa6a6542e344.
\end{center}
Consequently $c_*\in K_1$.  The separator for the old two-frame cone
does not separate this enlarged cone.

For clarity, its support inequality in the eight observable coordinates is
$\ell_{c_*}\geq0$, where for a general $c$
\begin{align}
\ell_c={}&d\left[1-\frac{E_{00}-E_{01}+E_{10}+E_{11}}2\right]
 +m(1+a_1-b_1-E_{11})\notag\\
 &+C(1-a_1+b_0-E_{10})+A(1-a_0+b_0-E_{00})\notag\\
 &+D(1-a_1+b_1-E_{11})
 +\alpha(1-a_0-b_1+E_{01})\notag\\
 &+\beta(1+a_1-b_0-E_{10}).
 \label{eq:visible}
\end{align}
The exact nine affine coefficients are included in the selector data.

Here is the full physical restoration of this inequality.  On a real
two-qubit block choose
\[
 A_0=B_0=\sigma_z,\qquad
 A_1=\frac{(1-u^2)\sigma_z+2u\sigma_x}{1+u^2},\qquad
 B_1=\frac{(1-v^2)\sigma_z+2v\sigma_x}{1+v^2}.
\]
Replace the observable coordinates in \eqref{eq:visible} by the
corresponding operators, denoting the result by $\ell_c(A,B)$.
In the computational basis $\psi=(W,X,Y,Z)^T$, direct substitution gives
\begin{equation}
 (1+u^2)(1+v^2)\psi^T\ell_c(A,B)\psi=F_c(u,v;\psi).
 \label{eq:physicalrestoration}
\end{equation}
For example, the six positive-event amplitudes are proportional to
$vW-X+uvY-uZ$, $uW-Y$, $Y$,
$uW+uvX-Y-vZ$, $vY-Z$, and $X+uZ$, respectively.
Thus \eqref{eq:gram}, with its nonnegative weights on the explicit testing
domain $u,v\geq0$, proves the operator inequality on every such block.
The sign gauges and continuity cover all angles and endpoints.  Real
operator positivity implies complex operator positivity by separating
the real and imaginary parts of a state vector.

Finally, two projective binary observables on one party decompose into
invariant subspaces of dimension at most two: if $U=A_0A_1$, then
$A_0UA_0=U^*$, so the spectral subspaces of $U$ pair at conjugate
eigenvalues.  On each resulting subspace the displayed Pauli form
applies, including the one-dimensional endpoints.  Applying this on
both parties proves $\ell_c\geq0$ for all quantum behaviors.  Binary
POVMs are covered by a common local projective dilation.  Thus the
source certificates here restore full-Q222 Bell inequalities.
The positive-coordinate requirements specify the native source chart;
validity follows from the restored full-parent identity, rather than
from those coordinate signs alone.

There is also a full-dimensional cone with a linear auxiliary selector.
Let $e_i$, $i=1,\ldots,6$, be the unit event rays and $y_i$ their exact
new-frame auxiliary vectors.  For
\begin{equation}
 \mathcal W=\{c:d\geq0,\quad c_i-d(c_*)_i\geq0,
                           \ i=1,\ldots,6\},
 \label{eq:chamber}
\end{equation}
set
\begin{equation}
 y(c)=d y_*+\sum_{i=1}^6(c_i-d(c_*)_i)y_i.
 \label{eq:chamberselector}
\end{equation}
The corresponding matrices are the same nonnegative combination of the
seven verified PSD matrices, proving $\mathcal W\subseteq K_1$.
The exact $302\times7$ selector is supplied in
\path{C1_REPAIRED_CHAMBER_AUXILIARY_MAP.json}.

\section{A rational selector on an open contact domain}

Put $\theta=(p,q,a,b,t)$, $\phi=(a,b,t,1)^T$, and require
$p,q>0$ and $abt\ne0$.
Define the $6\times7$ polynomial contact matrix $\mathcal J(\theta)$
as follows: its first four rows give $\mathcal P_c(p,q)\phi/4$, and
its last two rows give
$\phi^T\mathcal P_{c,u}(p,q)\phi/4$ and
$\phi^T\mathcal P_{c,v}(p,q)\phi/4$.
These expressions are linear in $c$.
Writing $\mathcal J=(\mathcal J_d\mid\mathcal J_e)$, set
\begin{equation}
 c(\theta)=\bigl(1,-\mathcal J_e(\theta)^{-1}
                         \mathcal J_d(\theta)\bigr),
 \qquad
 \Omega=\{p,q>0, abt\ne0, \det\mathcal J_e\ne0,
                          c_1,\ldots,c_6>0\}.
 \label{eq:sourcechart}
\end{equation}
This is a rational definition from the entire parent.  The data print
all 42 entries of $\mathcal J$; each was independently rederived from
\eqref{eq:parent}.  At the source in \eqref{eq:seed},
\begin{equation}
 \det\mathcal J_e
   =-\frac{7033805705872535849475830606568681}
                  {344646229622500000000}\ne0.
 \label{eq:minor}
\end{equation}

Make the source-only change of variables
\begin{equation}
 u=pU,\qquad v=qV,\qquad \psi=\diag(\phi)\chi.
 \label{eq:normalize}
\end{equation}
In sector $g$, let $D_g=\diag(p^rq^s\phi_i)$ on the retained features.
The normalized Gram is
$\widetilde G_g=p^{\alpha_g}q^{\beta_g}D_gG_gD_g$.
Its contact feature vector is $\one$.
For $S_g=[I_{n_g-1}\mid-\one]$, the matrices with this kernel are exactly
\begin{equation}
 \widetilde G_g=S_g^TH_gS_g,
 \qquad H_g\in\mathbb S^{n_g-1},
 \qquad H_g\succeq0\ \Longleftrightarrow\
                              \widetilde G_g\succeq0.
 \label{eq:compress}
\end{equation}
Here $H_g$ is the leading principal submatrix of $\widetilde G_g$.
The equivalence follows also directly from the surjectivity of $S_g$.

Let $h\in\R^{392}$ list the upper entries of the four $H_g$, in
lexicographic order $(g,i,j)$.  Let $\mathfrak M$ be the 148 scalar
monomials $U^rV^s\chi_i\chi_j$ present in \eqref{eq:gram}, in
lexicographic order $(r,s,i,j)$, $i\leq j$.
Denote by $E$ the $148\times7$ source coefficient matrix of
$(1+u)(1+v)F_c$.  The normalized target is
\begin{equation}
 b(\theta)_{r,s,i,j}
     =p^rq^s\phi_i\phi_j(Ec(\theta))_{r,s,i,j}.
 \label{eq:b}
\end{equation}
The difference features $(V_{g,i}-V_{g,n_g})$ define a constant rational
coefficient matrix $R\in\Q^{148\times392}$ such that the entire matching
equation is $Rh=b(\theta)$.

Here are the exact linear identities needed below.  Let $C$ have columns
\begin{equation}
 C_{:,r,s,i,j}=
  (\delta_{0i}+\delta_{0j},\ldots,
        \delta_{3i}+\delta_{3j},r,s)^T.
 \label{eq:C}
\end{equation}
Then
\begin{equation}
 CR=0,\quad \operatorname{rank}C=6,\quad
 \operatorname{rank}R=142,\quad \operatorname{im}R=\ker C.
 \label{eq:rank}
\end{equation}
The first identity follows because a product of two difference features
has zero value and first derivatives at the contact.  For an explicit
certificate of the rank assertion, choose the lexicographically first
142 independent rows $I$ of $R$ and the first 142 independent columns
$J$ of $R_I$.  The supplied rational matrix
\begin{equation}
 B=R_{I,J},\qquad P=B^{-1}
 \label{eq:inverse}
\end{equation}
obeys $BP=I_{142}$.  Its 1387 nonzero entries and the complete coefficient
maps are in
\path{C1_ENLARGED_REGULAR_CONTACT_SYSTEM.json}.
Thus \eqref{eq:rank} is certified by explicit finite matrix identities.
Since the complete contact equations vanish,
\begin{equation}
 Cb(\theta)=0\qquad(\theta\in\Omega).
 \label{eq:compatible}
\end{equation}
The multiplier derivatives contribute zero because the parent value is
zero at the contact.  In particular no state equation has been omitted.

Let $h_*$ be obtained from \eqref{eq:seedfact} by
\eqref{eq:normalize}--\eqref{eq:compress}, let $b_*=b(\theta_*)$, and
let $J^c$ be the remaining 250 columns.  Define the rational selector
\begin{equation}
 h_{J^c}(\theta)=(h_*)_{J^c},\qquad
 h_J(\theta)=(h_*)_J+P\bigl(b(\theta)_I-(b_*)_I\bigr).
 \label{eq:selector}
\end{equation}
Equations \eqref{eq:rank}--\eqref{eq:compatible} imply $Rh=b(\theta)$:
the difference has zero $I$ rows and lies in the 142-dimensional space
$\operatorname{im}R$, on which those rows are injective.

\begin{theorem}[Explicit open domain]
Let
\begin{equation}
 \mathcal U=\left\{\theta\in\Omega:
       \|b(\theta)-b_*\|_\infty<\frac1{3284140032}\right\}.
 \label{eq:U}
\end{equation}
For every $\theta\in\mathcal U$, \eqref{eq:selector} gives four positive
definite compressed Grams and hence an exact certificate
$c(\theta)\in K_1$.  The set $\mathcal U$ contains an open neighborhood
of $\theta_*$, and its image contains a five-dimensional family of
distinct stationary sources.
\end{theorem}

\begin{proof}
The first $n_g-1$ rows of the factors in \eqref{eq:seedfact} are invertible.
The source normalizations are invertible too, so every $(H_g)_*$ is
positive definite.  Exact rational inverse traces give the bounds
\[
\begin{array}{c|rrrr}
g&0&1&2&3\\\hline
\tr((H_g)_*^{-1})\leq&2^{12}&2^{17}&2^{22}&2^{17}\\
L_g&364&437/2&409/2&783/2
\end{array}
\]
The exact inverse traces are supplied with the selector data.
Since $\lambda_{\min}(H)\geq1/\tr(H^{-1})$ for $H\succ0$,
\begin{equation}
 (H_g)_*\succeq2^{-22}I.
 \label{eq:margin}
\end{equation}
For completeness, the displayed row bounds use the following explicit
definition.  For a pivot entry indexed by $(g,i,j)$, put
$a_{gij}=\sum_{k=1}^{142}|P_{\operatorname{pos}(g,i,j),k}|$; put
$a_{gij}=0$ for a fixed free entry and reflect it symmetrically.
Then $L_g=\max_i\sum_j a_{gij}$.
Thus \eqref{eq:selector} gives
\[
 \|H_g(\theta)-(H_g)_*\|_2
 \leq\|H_g(\theta)-(H_g)_*\|_\infty
 \leq\frac{783}{2}\|b(\theta)-b_*\|_\infty<2^{-23}.
\]
Together with \eqref{eq:margin}, this proves
$H_g(\theta)\succ2^{-23}I$.

Restore the original Grams by the explicit rational formula
\begin{equation}
 G_g(\theta)=p^{-\alpha_g}q^{-\beta_g}
 D_g^{-1}S_g^TH_g(\theta)S_gD_g^{-1}.
 \label{eq:restore}
\end{equation}
All denominators are nonzero on $\Omega$.  They depend only on the
source; the restored matrices satisfy the formal polynomial identity
\eqref{eq:gram} in independent $u,v,\psi$.
The original auxiliary $y_j(\theta)$ is the entry of $G_g(\theta)$ at
the $j$th listed free position of the literal source pencil.

The inequalities defining $\Omega$ are strict at \eqref{eq:seed} and
the target is continuous there, so $\mathcal U$ contains an open
neighborhood.  Differentiating the six contact equations gives
\[
 \partial_{\theta_k}c_e
 =-\mathcal J_e^{-1}(\partial_{\theta_k}\mathcal J)c.
\]
At the seed this $6\times5$ rational matrix has rank five.  This proves
the asserted local dimension without identifying testing variables with
source variables.
\end{proof}

\section{A short rational source family and the selector's limitation}

Fix the four contact coordinates
$(p,q,b,t)=(1/13,16,-1/32,1/50)$ and vary $a=\sigma$.
Put $Q_\sigma=83200\sigma^2+17476053\sigma+119$.  Solving the entire
contact system gives $d=1$ and
\begin{align}
 m(\sigma)&=\frac{3499569000}
 {(166400\sigma-219)Q_\sigma},\notag\\
 C(\sigma)&=\frac{177021000(2\sigma-1)}
 {(50\sigma-13)Q_\sigma},\notag\\
 A(\sigma)&=\frac{2600(12800\sigma^2+7122\sigma-5219)}{Q_\sigma},\notag\\
 D(\sigma)&=-\frac{2080000\sigma^3+2680525\sigma^2
                         -218723632\sigma-621061}
                    {2(25\sigma-5219)Q_\sigma},\notag\\
 \alpha(\sigma)&=\frac{3204(83200\sigma^2+53\sigma+119)}{17Q_\sigma},\notag\\
 \beta(\sigma)&=-\frac{13(121600\sigma^2-17362101\sigma-83623)}
                         {38Q_\sigma}.
 \label{eq:family}
\end{align}
These formulas are exact stationary-source identities, including all
four state equations and both angular equations.
Their contact minor is
\begin{equation}
 \det\mathcal J_e(\sigma)=
 \frac{6165476003(25\sigma-5219)(50\sigma-13)
                       (166400\sigma-219)Q_\sigma}
      {344646229622500000000},
 \label{eq:familyminor}
\end{equation}
which is nonzero throughout $139\leq\sigma\leq141$.

All six coefficients are positive for $139\leq\sigma\leq141$.
Here is a direct uniform sign verification.  Write each displayed
fraction as $N_i(s)/D_i(s)$ after $\sigma=140+s$ and expand the
numerator and denominator, without cancelling their displayed integer
factors.  Each constant has greater absolute value than the sum of
the remaining coefficient magnitudes.  The positive differences are
\[
\begin{array}{c|r|r}
 i& |N_i(0)|-\sum_{k\geq1}|N_{ik}|
    &|D_i(0)|-\sum_{k\geq1}|D_{ik}|\\\hline
 m&3499569000&93348433945876566\\
 C&49034817000&27997200202682\\
 A&645496641400&4036512286\\
 D&5605722176566&13952538819296\\
 \alpha&5149943908344&68620708862\\
 \beta&397933432&153387466868
\end{array}
\]
The constant numerator and denominator have the same sign in every row.
This proves the assertion on $|s|\leq1$.

For an explicit interval on which the selector is guaranteed to work,
write each target entry in \eqref{eq:b} as a reduced fraction
$N(s)/D(s)$.  Let $\|N\|_1=\sum|N_k|$ and
$\|N'\|_1=\sum k|N_k|$, and set
$d_0=|D_0|-\sum_{k\geq1}|D_k|>0$.
For $|s|\leq1$,
\[
 \left|\frac{\dd}{\dd s}\frac{N(s)}{D(s)}\right|
 \leq\frac{\|N'\|_1\|D\|_1+\|N\|_1\|D'\|_1}{d_0^2}.
\]
Taking the maximum over the explicit target entries gives the exact bound
\[
 B=\frac{18674611204409407977433140269829677208000}
          {37399703413558479094297615998554952613}<500.
\]
Therefore every rational or real $\sigma$ satisfying
\begin{equation}
 |\sigma-140|\leq2^{-42}
 \label{eq:familyinterval}
\end{equation}
has $\|b(\sigma)-b_*\|_\infty\leq B2^{-42}
<1/(2\cdot3284140032)$ and is certified by
\eqref{eq:selector}--\eqref{eq:restore}.

The same fixed choice of 250 free entries cannot be asserted on the
whole regular contact domain.  For example, \eqref{eq:family} at
$\sigma=70$ has six strictly positive coefficients and a nonsingular
contact matrix.  The restored compressed selector has
\begin{equation}
 (H_2(70))_{00}
 =-\frac{4276129302757499074690240852689277812474790802751901913517}
 {9456495976921764179751029215348840318787988500000000000}<-400.
 \label{eq:selectorfailure}
\end{equation}
It is thus not PSD.  Even at $\sigma=141$, the first nonpositive
Schur pivot in $H_0$ is its sixth pivot and is strictly negative;
the exact rational value is supplied.  These calculations refute the
universal use of \emph{this selector}.  They do not refute membership in
$K_1$: other values of the 250 free coordinates remain available.

\section{The precise remaining adequacy question}

The general C1 contact problem has an explicit finite form.  For an
arbitrary regular contact, let $h_{J^c}=z\in\R^{250}$ and set
\begin{equation}
 h_J=P\bigl(b_I-R_{I,J^c}z\bigr).
 \label{eq:general}
\end{equation}
Its four matrices have orders $(13,14,13,14)$.  The other canonical
orientation is the rational contact map
\begin{equation}
 \tau(p,q,a,b,t)=\left(p^{-1},q,
       \frac{qa-b}{t+q},\frac{-a-qb}{t+q},\frac{1-qt}{t+q}\right),
 \label{eq:reflection}
\end{equation}
with source permutation
$\sigma c=(d,C,m,\alpha,\beta,A,D)$.
The exact parent congruence behind this map is
\[
 u^2(1+v^2)\mathcal P_{\sigma c}(1/u,v)
     =G(v)\mathcal P_c(u,v)G(v)^T,
\quad
G(v)=\begin{pmatrix}v&-1&0&0\\-1&-v&0&0\\
0&0&-v&1\\0&0&1&v\end{pmatrix}.
\]
On the regular two-orientation domain, the universal statement still
needed is that for every contact there is an orientation and a value
of $z$ for which all four matrices in \eqref{eq:general} are PSD.
The construction above proves this statement on $\mathcal U$ with one
explicit value of $z$.  It supplies no uniform selector on the
complement.  Completing C1 would also leave the independent whole-C2
converse needed by the order-13744 compiler.

In particular, a positive Gram margin at one contact establishes an
open family, but it neither gives a finite cover of escaping or singular
contacts nor licenses a constant free-coordinate choice everywhere.
Equation \eqref{eq:selectorfailure} identifies this limitation exactly.
No conclusion about C160 reverse inclusion is inferred from the C1
certificate.

\section*{Exact data and reproduction}

The directory \path{c1_contact_selector_data/} contains the literal
source pencil, the repaired source's full auxiliary vector and four
positive factorizations, the event-ray selector, the complete constant
contact coefficient maps, and the newly evaluated open-domain and
one-parameter selectors.  The independent exact reproduction script
\path{check_c1_contact_selector.py} uses rational arithmetic and
SymPy, and invokes no numerical optimizer.  It directly checks the
current source-array substitution, all seven full-parent source
columns, all six contact equations, the difference-feature coefficient
map, its inverse-minor certificate, the strict Gram bounds, and the
stated rational family and failures.  The script and data also work
without the historical archive when kept together.

\end{document}
